\cleardoublepage
\chapternonum{致谢}
\ifthenelse{\equal{\BlindReview}{true}}
{% For blind review
}
{% For submission
    我还清楚地记得，八年前第一次走进求是园的那个夏天。溽热扑面而来，满眼的草地和绿植绿得发亮，像是能滴出水来。那时的我，怀揣着对学术最朴素的热爱和对未来最饱满的期待，觉得一切皆有可能。八年过去，求是园的草木枯荣了八回，我也在这座园子里完成了人生中最重要的一次蜕变。如今论文付梓，我深知，这段“痛苦但值得”的旅程，从来不是一个人走完的。所有那些在我迷惘时指过路、在我跌倒时伸过手的人，都值得被郑重地写在这里。

    最该感谢的是我的指导老师黄劲教授。他既是我的科研指导老师，更是我的人生引路人。感谢他在九年前的那个暑假接受了跨专业的我。黄老师对学术的热爱和追求、严谨的态度，让我非常敬佩，也为我培养了正确的学术观。黄老师更是一位出色的老师，善于帮我观察和总结身上的缺点，尽力帮我改正，就如他所说，他在“debug”一个人。黄老师更是我遇见过的最认真负责的老师，早年间我的代码工程功底不扎实的时候，他曾经陪我调试代码到深夜凌晨两点，我们也经常讨论学术问题到深夜。在论文投稿上，黄老师事必躬亲，从选题到实验设计到论文撰写，主动提供大量的帮助。同时黄老师非常支持我的想法，当我有一些点子的时候，他会非常开心地支持；当我们想法不一致时，他也会客观地、耐心地就事论事和我讨论分析。此外，黄老师始终鼓励我坚持下去，尽管我内心打了无数次退堂鼓，甚至提出过转硕。没有这份鼓励，我是无法完成这段博士学业的。黄老师专业的学术能力为我提供了尤其宝贵的指导，他分析问题、解决问题的严谨思路，更是帮助我从一个科研新手一步步走上了科学化、规范化的科研道路。除去以上的指导值得我终生铭记，黄老师对人对世界的看法也深深影响了我。黄老师对人对事的分析看法角度独到，帮助我更好地理解了这个世界。总结起来，就是学会认识自己，学会认识世界：知道自己的过度完美主义、不严谨不逻辑、拖延症、粗枝大叶，也要知道自己的不畏挑战、喜欢深挖、能够坚持的特性；知道这个世界会有很多不如意，会有丰富多彩的评价体系，会在你每次有希望的时候给你当头一棒，也会在你绝望的时候给你曙光。而最重要的，就是既坚持一些真正对的、有意义、有价值的事，同时改变自己，适应这个世界。综上所述，认识黄老师并且在他的指导下获得博士学位，是我这一生最宝贵的经历。

    感谢在读博期间认识的所有同学和师兄弟。感谢陈炯师兄对我的宝贵指导，始终耐心地、严谨地回答我的问题，并且有求必应。感谢王天宇师兄的指导和陪伴，给我树立了一个热爱学术的榜样，让我明白热爱学术不应该被身份、职业所限制。感谢王诗怡师姐在生活中的陪伴，给博士生涯增添了欢乐。感谢郑一村师兄、林敏良师兄在技术、生活、规划上的帮助。感谢林锦铿师弟多年的帮助和合作，坐在我后面的他几乎陪我走过了整个博士旅程，是我成长路上的重要帮手。感谢孙浩然师弟和张小虎师弟在多方面的帮助和鼓励。感谢王润冬师弟和秦际州师弟在论文实验上的帮助以及生活中的陪伴。额外感谢尹君州、张峻溪在论文合作中提供的帮助。感谢石泽云、方贤忠、蒋静、高宏玉、郑驭聪、孙典圣、李铭扬、杨博淳、章沈柯、吴双、靳奕扬、徐翊钧、俞枳枰在组里的帮助。

    感谢博士期间认识的两位好友，江彦开和冯旭东。江彦开是我博士期间最重要的好伙伴，既有一起快乐玩耍的时光，也在我低谷的时候给予我大量的帮助，更在学术、职业发展、就业选择上给了我大量好的建议和帮助。冯旭东和我工作方向相似，起初我们在物理仿真领域时常一起探讨，后来经常一起讨论各类技术发展、社会发展和人生理想。他待人真诚，是不可多得的知心好友，同样在职业发展上给予了我重要的帮助。

    感谢我的父母，始终鼓励和支持我。漫漫三十年，是他们从小到大的抚养让我能够度过漫长的学习生涯，完成这份凝聚了我的心血、也是他们的心血的博士论文。在生活上他们给予了我重要的帮助，给了我优渥的条件可以专注学业。我的任何决定他们都无条件地支持，他们是我最最坚强的后盾。

    八年求是园，一生求是情。感谢这段时光，感谢所有遇见。未来的路还很长，我会带着这份感恩与成长，踏实走好每一步。愿不负所学，不负所爱，不负自己。

    \begin{flushright}
        \vspace*{\fill}
        \StudentName

        2026年5月于杭州
    \end{flushright}
}
